\section{An explicit odd-dimensional nonselfdual orthogonal bound}
\label{sec:orthogonal-odd-nonselfdual}

\paragraph{Status and scope.}
This is a separate new author derivation.  The complete written
proof and exact supplement pass the author check; whole-proof
reviews by /root and /root/g34\_eightfolds have also passed, with
a separate exact peer supplement.  It leaves D-OL-06--11 unchanged.  All modules
below are ordinary complex modules of the natural orthogonal
derived group, rather than arbitrary spin modules.  The field is
of odd order.  All compatible finite-projective-order normalizers
are included, with arbitrary field components.

Put
\[
 \delta_*:=\frac{1-(19/20)^9}{2},\qquad
             \frac16<\delta_*<\frac3{16},\qquad
             \frac6{12\delta_*-2}<28.                       \tag{OO.1}
\]
The three displayed inequalities are exact rational inequalities.

\begin{theorem}[D-OL-12]
\label{D-OL-12}
Let $q\ge5$ be odd and let $V$ be a nondegenerate quadratic
space of odd dimension $d\ge1135$ over $\mathbf F_q$.
Write $S=\Omega(V)$ and $O=O(V)$.  Let $\rho$ be a nonselfdual
ordinary irreducible complex $S$-module of degree $a$.
Every nonscalar compatible normalizer $T$ of finite projective
order satisfies
\[
                         \nu(T)>\delta_*a.                  \tag{OO.2}
\]
An irreducibly acting normal central product of such modules
cannot supply period-image rank nine for its actual six-dimensional
multiplicity variation.  Mixed odd fields and odd natural
dimensions, represented central quotients, scalars and all
compatible factor permutations are allowed.  No determinant
assumption or selfdual structure on the inducing representation
is needed for this assertion.
\end{theorem}

\begin{proof}
We first prove the inner support bound, without using
nonselfduality.  The natural group $S$ has trivial center in odd
dimension and is simple in this range.  Its projective kernel
on any nontrivial irreducible complex module is therefore trivial:
that kernel is normal, and if it were all of $S$ the perfect
group $S$ would act by scalars.  It suffices to take a nonidentity
prime-order power $g$ of an inner operator, because taking powers
cannot increase $\nu$.

Suppose first that one of the natural eigenspaces of $g$ for
$\epsilon\in\{1,-1\}$ contains a hyperbolic pair $e,f$.
Write
\[
 V=Fe\oplus W\oplus Ff,\qquad
 Q(xe+w+yf)=xy+Q_W(w),\qquad N=d-2.                          \tag{OO.3}
\]
The two opposite abelian line radicals $U_e,U_f$ of D-OL-06
are conjugate inside $S$ and generate $S$.  Thus
\[
                       b:=\dim\rho^{U_e}\le a/2.           \tag{OO.4}
\]
The formula for $u_e(v)$ in (OS.3) gives the exact identity
\[
              gu_e(v)g^{-1}=u_e(Av),\qquad A=\epsilon g|_W.
                                                               \tag{OO.5}
\]
The same isometry $A$ acts on the dual weight vectors, up to
inversion, which does not alter fixed points or orbit lengths.
If $\epsilon=-1$, then $g$ has order two and $A^2=I$.
If $\epsilon=1$, $A$ has the prime order of $g$.  In either
case $A$ is nonidentity: otherwise $g=\epsilon I$ on $V$,
and neither the identity nor $-I$ is the nonidentity element
of $S$ under consideration.  Consequently the action on $W$
has a prime order $l$ and moving rank $r\ge1$.

It is important here to allow $r=1$.  For example, an element
which is minus a reflection can lie in $S$ in odd dimension;
the replacement $A=\epsilon g|_W$ is then a reflection.  No
rank-two lower bound is imposed in this argument.

Let $X=\ker(A-I)$.  It has dimension $N-r$, and its restricted
polar radical has dimension at most $r$.  For $r\le N/2$, the
quadratic Gauss error (OS.10) is at most
\[
              q^{(N-r)-(N-2r)/2}=q^{N/2}.
\]
For $r>N/2$, the trivial bound on the Gauss sum is at most
$q^{N-r}\le q^{N/2}$.  As in D-OL-06, the nonzero Levi orbits
are exactly the norm levels
$\mathcal O_t=\{x\ne0:Q_W(x)=t\}$, and weight multiplicities
are constant on every such orbit.  Hence the proportion fixed
by $A$ in any one of them is at most
\[
 \begin{split}
 \frac{q^{N-r-1}+q^{\lceil N/2\rceil}}
      {q^{N-1}-q^{\lceil N/2\rceil}-1}
 &\le \frac{q^{-1}+q^{-9}}{1-q^{-9}-q^{-19}}\\
 &\le \frac{5^{-1}+5^{-9}}{1-5^{-9}-5^{-19}}<\frac14.
 \end{split}                                                 \tag{OO.6}
\]
Here $N\ge20$; dividing by $q^{N-1}$ gives the exponent
$1-\lfloor N/2\rfloor\le-9$.  The numerator decreases and
the positive denominator increases with $q$, so the final
inequality is uniform, not a finite-field extrapolation.
Every remaining weight orbit has length $l$.  The component
recovery argument (OS.14), valid for arbitrary scalar phases,
now yields
\[
           \nu(\rho(g))\ge(1-1/l)\frac34(a-b)
                           \ge\frac3{16}a>\delta_*a.        \tag{OO.7}
\]

Suppose next that neither natural $\epsilon$-eigenspace contains
a hyperbolic pair.  If the prime order of $g$ differs from the
characteristic, $g$ is semisimple.  Its $1$- and $-1$-eigenspaces
are nondegenerate: their orthogonal complements are the other
primary spaces.  A nondegenerate quadratic space of dimension
at least three over a finite field of odd order contains a
hyperbolic pair.  Thus both eigenspaces have dimension at most
two, and
\[
            \min_{\epsilon=\pm1}\operatorname{rank}
                      (g-\epsilon I)\ge d-2.               \tag{OO.8}
\]
This also shows directly that the present case cannot occur for
an involution when $d>4$.

If $g$ has the characteristic prime order, it is unipotent.
Put $r=\operatorname{rank}(g-I)$.  The fixed space has dimension
$d-r$, and its polar radical has dimension at most $r$, because
its perpendicular space is $\operatorname{im}(g-I)$.
Its polar rank is at least $d-2r$.  Polar rank at least three
would give a nondegenerate three-dimensional subspace and hence
a fixed hyperbolic pair.  Therefore
\[
          r\ge(d-2)/2,\qquad \operatorname{rank}(g+I)=d.
                                                               \tag{OO.9}
\]
Since $d\ge1135$ is odd, the integer lower bound in (OO.9)
is at least $567$.  Equations (OO.8)--(OO.9) prove that the
natural support restricted to $\mu_2=\{1,-1\}$ is at least
$567$ in every remaining case.

Inflate $\rho$ along the ordinary simply connected finite cover
$\operatorname{Spin}(V)\longrightarrow S$.  It remains nontrivial
and irreducible, and has the same natural orthogonal action for
the support calculation.  The cover is nonexceptional in these
large dimensions; the precise cover scope is the one used in
D-OL-10.  Apply
\cite[Theorem~4.3.6, equation~(4.3.2), pp.~1917--1918;
complete proof pp.~1918--1923; Remark~4.3.7, p.~1923]{HOS:LST2011}
with the support threshold $567=7\cdot9^2$.
The remark permits exactly the restricted scalar set $\mu_2$
for the spin group.  The published constant is $19/20$ for
$2\le q<43$, and $1/(\sqrt q-1)\le19/20$ for the remaining
fields.  No faithful-character hypothesis is required.
Consequently
\[
             |\chi(g)|/a<(19/20)^9,
             \qquad \nu(\rho(g))>\delta_*a.                 \tag{OO.10}
\]
The second inequality follows from
$|\chi(g)|\ge a-2\nu(\rho(g))$.  This proves the inner
bound for every nonidentity $g\in S$.

We now use nonselfduality to determine the compatible linear
normalizers.  In odd dimension a linear similitude with multiplier
$\mu$ has $(\det h)^2=\mu^d$, so $\mu$ is a square.  Hence
\[
 GO(V)=F^*O(V),\qquad O(V)=\langle-I\rangle\times SO(V),
                    \qquad SO(V)/S\simeq C_2.               \tag{OO.11}
\]
The last quotient is the determinant/spinor-norm description
\cite[Section~2, pp.~41--43]{OL:KTV2023}.  Thus there is only
one nontrivial linear outer coset after natural scalars and
inner automorphisms have been removed.

If that coset stabilized $\rho$, then all of $O$ would stabilize
its isomorphism class: the additional factor $-I$ is central.
Take an irreducible constituent of $\operatorname{Ind}_S^O\rho$.
Clifford theory makes its restriction a positive multiple of
$\rho$, since the conjugacy orbit of $\rho$ would be a singleton.
But all irreducible characters of $O$ are real-valued, by the
published full-group reality theorem
\cite[Theorem~1, p.~44; attribution and exact domain
pp.~39--40]{OL:KTV2023}.  Its restricted character is therefore
real, forcing $\rho$ to be selfdual, a contradiction.
Every compatible linear normalizer is consequently a scalar
times an inner operator, and satisfies (OO.10).

For completeness, this argument is applied to every linear
component of the full natural semilinear-normalizer description,
not to a chosen subgroup of automorphisms.  The description and
its cardinality/centralizer proof in D-OL-07 apply for arbitrary
field degree in these dimensions; the complete published outer
table is \cite[Section~2.2.2, Table~2.1, p.~14]{HES:LiXia2022}.
For a general compatible $T$, take a nonscalar power of prime
projective order $l$.  If its field component is trivial, the
preceding paragraph applies.  Otherwise $q=q_0^l$.  The full
norm and quadratic-form descent (ON.7)--(ON.9) applies to its
possibly mixed linear--field lift: after scalar and form
rescaling, that lift has order $l$, fixes an $F$-basis, and its
quadratic form is defined over $F_{q_0}$.  A hyperbolic pair over
$F_{q_0}$ gives a line radical on which it acts by Frobenius.
The fixed-weight fraction is less than $1/4$, uniformly in
the arbitrary prime field-extension degree.  Equations
(ON.7)--(ON.9) therefore give
\[
                         \nu(T)\ge3a/16>\delta_*a.          \tag{OO.12}
\]
There is no assumption that the original lift was pure Frobenius.
Taking the power only decreased projective support, so (OO.2)
holds for the original normalizer.

On tensor products, convolution retains the support fraction
of any active nonscalar factor.  A nontrivial factor permutation
has projective support at least one quarter of the full degree,
which is greater than $\delta_*$.  A nonidentity scalar has full
actual moving rank.  Thus every genuine inertia moves more than
$\delta_*e$ dimensions.

The ordinary factor degree estimate
\cite[Theorem~3.7, pp.~165--167]{OL:GLT2024} gives
$a_i\ge q_i^{d_i-4}(q_i-1)>28$, so $e>28$.
The actual intermediate-cover Euler identity
$6=2e(h-1)+D$ excludes genus at least two.  In genus one the
moving sum is six, hence there are at most six genuine marks
and the actual period-image rank is at most six.  Rank nine
therefore requires genus zero and at least twelve genuine
inertia points.  It would force
\[
             2e+6>12\delta_*e,
             \qquad e<\frac6{12\delta_*-2}<28,              \tag{OO.13}
\]
a contradiction.  The argument uses the actual intermediate
variation and does not identify the inducing character field
with the original imaginary quadratic field.
\end{proof}

\begin{corollary}[D-OL-13]
\label{D-OL-13}
In the main imaginary-quadratic sixfold problem, a transitive
irreducibly acting normal central product of ordinary orthogonal
factors of odd natural dimension over an odd field $q\ge5$
must have
\[
                              d\le1133.                    \tag{OO.14}
\]
If the factor module is selfdual, the sharper earlier bound
$d\le21$ applies.  These are necessary bounds, not existence
claims for the remaining parameters.
\end{corollary}

\begin{proof}
Selfduality is preserved by automorphic pullback, so transitivity
makes all factor modules selfdual or all nonselfdual.  The first
case is D-OL-09.  The second is D-OL-12.  Since $d$ is odd,
the exclusion $d\ge1135$ is equivalent to (OO.14).
\end{proof}

\paragraph{Why this calculation does not include $q=3$.}
The rank-one possibility in (OO.5) cannot be discarded.  For
$Q_W=\sum_{j=1}^s x_jy_j+z^2$ and $A(x,y,z)=(x,y,-z)$,
the fixed fraction on the nonzero singular level is
\[
 \frac{q^{2s-1}+(q-1)q^{s-1}-1}{q^{2s}-1}.                  \tag{OO.15}
\]
The numerator and denominator follow by successively adding
hyperbolic planes: $xy=t$ has $2q-1$ solutions for $t=0$ and
$q-1$ for $t\ne0$, giving the displayed norm-level counts by
induction, with the zero vector removed in both counts.
At $q=3$ this exceeds $1/3$ when $s>1$, so the uniform
$1/4$ bound used here is false in that case.  Minus a reflection
has spinor norm one on a natural space of dimension $4k+1$,
as is seen by writing its action on the $2k$ hyperbolic planes
as products of two reflections with norm product $-1$.
Thus the rank-one weight action is an actual inner possibility.
This is a limitation of this particular estimate and asserts
no existence of a characteristic-three source.

\paragraph{Source and certificate scope.}
The new steps are the signed hyperbolic-pair action, its rank-one
Gauss bound, the odd-dimensional nonselfdual normalizer argument,
and their exact source consequence.  The previously reviewed
radical generation, arbitrary-field descent, degree estimate,
and restricted-support character theorem are used with their
stated scopes.  The exact supplement checks all constants and
finite signed-action/orbit-count witnesses; the universal proof
is the one written above.
